% Isolated final-INCOMPLETE proposal B176. No frozen source is replaced.
\subsection*{The whole four-generator contact closure}
\statement{B.176}{Proved}{The three complete single-parameter contact families form a proper closed subgroup of the known contact kernel, separated from an explicit mixed contact by a finite quotient of the original ambient group}
\label{stmt:B176}
Use the actual filled group \(\Gamma _0\) and the marked source
\(H_5=\langle A_s,B_s,D_s,E_s,m\rangle\) of B.169--B.175.
Abbreviate its contact letters by \(A,B,D,E\), without changing any
of their original whiskers. Write \(\mathcal H_5\) for its actual
closure in \(\widehat{\Gamma _0}\). The filled projection and section give
\[
 p:\mathcal H_5\longrightarrow
 J_0=\widehat{\mathbb Z}_A\times\widehat{\mathbb Z}_m,
 \qquad
 \mathcal I_5=\ker p,
\]
with \(p(A)=A,p(m)=m\), and \(p(B)=p(D)=p(E)=1\).
All closures in the following formulas are actual ambient closures.
Put
\[
 \begin{gathered}
 J_4^{\rm ord}=\langle m,B,D,E\rangle,\qquad
 \mathcal J_4=\overline{J_4^{\rm ord}},\qquad
 \mathcal K_4=\mathcal J_4\cap\mathcal I_5,\\
 c=[m,D],\qquad w=[c,A],\qquad [g,h]=ghg^{-1}h^{-1}.
 \end{gathered}
\]
There is a finite quotient of the original \(\Gamma _0\) in which the
image of \(w\) lies outside the entire image of \(J_4^{\rm ord}\).
Consequently
\begin{equation}\label{eq:B176-kernel}
\begin{gathered}
 w\notin\mathcal J_4,\qquad
 \mathcal J_4=\mathcal K_4\rtimes\widehat{\mathbb Z}_m,\\
 \mathcal K_4
 =\overline{\langle\!\langle B,D,E\rangle\!\rangle}_{\mathcal J_4}
 =\overline{\langle m^z C m^{-z}:
              z\in\widehat{\mathbb Z},\ C\in\{B,D,E\}\rangle}
 \subsetneq\mathcal I_5.
\end{gathered}
\end{equation}
The ordinary kernel of \(J_4^{\rm ord}\to\mathbb Z_m\) is generated by
the corresponding integer-indexed conjugates; its actual closure is
\(\mathcal K_4\). Furthermore
\[
 c\in\mathcal K_4,\qquad AcA^{-1}\in\mathcal I_5\setminus\mathcal J_4.
\]
Thus \(A\) normalizes neither \(\mathcal J_4\) nor \(\mathcal K_4\).

Let
\[
 \pi:\widehat P_0=
 (\widehat{\mathbb Z}_A\times\widehat{\mathbb Z}_m)
       \amalg\widehat{F(B,D,E)}\longrightarrow\mathcal H_5,
 \qquad C_0=\ker\pi,
\]
be the actual compact presentation map, and put
\(P_{4,c}=\overline{\langle m,B,D,E\rangle}\le\widehat P_0\).
Then the same marked words satisfy
\begin{equation}\label{eq:B176-preimage}
 w\notin C_0P_{4,c},\qquad AcA^{-1}\notin C_0P_{4,c}.
\end{equation}
For the injective connecting map \(\partial\) of B.164,
\begin{equation}\label{eq:B176-cohomology}
 \partial(\mathcal K_4)\subsetneq
 \partial(\mathcal I_5)\subseteq\partial(\mathcal I).
\end{equation}
These are compact closed subspaces of the actual nonabelian cohomology
space. The witness \(\partial(w)\) already belongs to the known source
image \(\partial(\mathcal I)\); it does not decide whether that image
is the entire required cohomology fiber.

\begin{proof}
\emph{The whole finite subgroup separator.}
Retain B.175's actual target
\[
 \begin{gathered}
 K_0=F(U,k,q,y)\rtimes S_*,\qquad S_*=F(m,V),\\
 K_1=K_0*_{S_*}(S_*\times\mathbb Z_{D_0}),\qquad
 G_\square=K_1*_{L}(L\times\mathbb Z_R),\\
 J=\langle S_*,U\rangle,\qquad
 H=\langle S_*,kUk^{-1}\rangle,\qquad
 L=\langle J,D_0HD_0^{-1}\rangle.
 \end{gathered}
\]
The source images are literally
\[
 A\mapsto R,\quad B\mapsto U^{-1},\quad
 D\mapsto UD_0U^{-1},\quad E\mapsto1,\quad m\mapsto m.
\]
Hence the whole \(J_4^{\rm ord}\) maps into the \(K_1\) vertex.
Writing \(U_1=m(U)\) and \(h=U_1^{-1}U\), the actual word \(c\) maps to
\[
                     U_1D_0hD_0^{-1}U^{-1}.
\]
B.175's finite fiber separator, characteristic refinement and finite
source action produce a finite head \(\overline K\) in which
\(\overline h\notin\overline H\). Its complete normal-form argument
in the intermediate finite-head HNN extension, followed by the finite
graph-core argument, gives a finite quotient \(Q_0\) of that extension
such that
\[
 c_0\notin L_0,
\]
where \(c_0,L_0\) are the images of this whole word and of \(L\).
This uses the full prefactors in \(c\), as in B.175. It is not a new
finite search or a claim about an intrinsic completion of \(L\).

Form the finite-vertex amalgam
\[
                   Z_0=Q_0*_{L_0}(L_0\times C_{2,R}).
\]
Its word \(w_0=[c_0,R]\) lies outside the whole \(Q_0\) vertex.
Indeed \(R\notin L_0\) in the second vertex and
\(c_0^{-1}\notin L_0\) in the first, so
\(Rc_0^{-1}R^{-1}\) is a reduced three-syllable word and is not in
\(Q_0\). If \([c_0,R]\) belonged to \(Q_0\), multiplication by
\(c_0^{-1}\) would contradict this conclusion. The argument is valid
with \(R\) of order two.

The natural map \(Z_0\to Q_0\times C_2\) is injective on each finite
vertex group. Its kernel acts freely on the ordinary amalgam tree and
is a finite-index finitely generated free group. This proves residual
finiteness directly: an element outside the kernel is separated by the
displayed finite map; a nontrivial element inside it is separated by a
finite quotient of the free group and then by its finite ambient core.
Since \(Q_0\) is finite, separate every nonidentity
\(w_0q_0^{-1}\), \(q_0\in Q_0\), by such a quotient and take the finite
product of these maps. The resulting finite homomorphism \(\nu\) has
\[
                         \nu(w_0)\notin\nu(Q_0).
\]

Compose the actual map from the index-two \(\Gamma _0^*\) with this
finite target. If \(K\) is its finite-index normal kernel, take
\(K^{\rm core}=\operatorname{core}_{\Gamma _0}(K)\le K\).
This is finite-index and normal in \(\Gamma _0\). Equality of the
images of \(w\) and any \(j\in J_4^{\rm ord}\) modulo \(K^{\rm core}\)
would imply equality in the preceding finite target, which is impossible.
Thus this is a finite quotient of the original ambient group separating
\(w\) from all of \(J_4^{\rm ord}\). Its continuous extension separates
\(w\) from \(\mathcal J_4\). No completion of an infinite edge group
has been identified with an ambient closure.

\emph{The actual kernel and its complete families.}
The filled marking gives \(p(\mathcal J_4)=\widehat{\mathbb Z}_m\)
with the literal \(m\)-section, so this is the full procyclic group.
Let \(N_4\) be the closed normal closure of \(B,D,E\) in
\(\mathcal J_4\). The quotient \(\mathcal J_4/N_4\) is topologically
generated by \(m\) and maps to \(\widehat{\mathbb Z}_m\) with that
generator sent to the standard generator. The canonical epimorphism
\(\widehat{\mathbb Z}\to\mathcal J_4/N_4\) has composite the identity
on \(\widehat{\mathbb Z}\), and is therefore an isomorphism. This
proves \(N_4=\mathcal K_4\) and the split formula.

Let \(L_4\) be the closed group generated by the entire family on the
right of \eqref{eq:B176-kernel}. It is invariant under every \(m^z\)
and contains \(B,D,E\). Those three elements normalize it internally.
They and \(m\) topologically generate \(\mathcal J_4\), so \(L_4\)
is closed and normal there. It is therefore precisely \(N_4\).
Integer powers are dense in each continuous procyclic power map; hence
it is also the closure of the subgroup generated by the integer-indexed
families. Ordinarily, killing \(B,D,E\) in \(J_4^{\rm ord}\) leaves
a cyclic group mapping isomorphically onto its actual \(\mathbb Z_m\)
image. Its kernel is the ordinary normal closure of those letters, which
is the group generated by their integer \(m\)-conjugates. This proves
its asserted actual-closure identity, without a full-topology assertion
for that possibly infinite-rank kernel.

The same argument using the actual commuting \(A,m\) section gives
\[
 \mathcal I_5=
 \overline{\langle\!\langle B,D,E\rangle\!\rangle}_{\mathcal H_5}
 =\overline{\langle A^i m^n C m^{-n}A^{-i}:
       i,n\in\widehat{\mathbb Z},\ C\in\{B,D,E\}\rangle}.
\]
After killing the three contacts, the quotient is generated by the
commuting \(A,m\) and maps identically to \(J_0\); the continuous
map from \(\widehat{\mathbb Z}^2\) is consequently an isomorphism.
The displayed family is invariant under both section generators and
contains the contacts, proving the equality of closed normal kernels.

Since \(c\in\mathcal K_4\), normality of \(\mathcal I_5\) in
\(\mathcal H_5\) gives \(w\in\mathcal I_5\). Its separation from
\(\mathcal J_4\) proves \(\mathcal K_4\subsetneq\mathcal I_5\).
The ordered identity
\[
                          AcA^{-1}=w^{-1}c
\]
puts \(AcA^{-1}\) in \(\mathcal I_5\) and outside \(\mathcal J_4\):
otherwise \(w\) would belong to \(\mathcal J_4\). This also proves
both nonnormality assertions.

\emph{The full source preimage and cohomology.}
In \(P_0=(\mathbb Z_A\times\mathbb Z_m)*F(B,D,E)\), the subgroup
\(P_4^{\rm ord}=\langle m,B,D,E\rangle\) is a free group of rank four.
The map \(A\mapsto1\), fixing its other four generators, is a retraction.
Its completion embeds as the compact retract \(P_{4,c}\) of
\(\widehat P_0\). Compactness and the marked generators imply
\(\pi(P_{4,c})=\mathcal J_4\). Since \(\pi\) is onto,
\[
                         \pi^{-1}(\mathcal J_4)=C_0P_{4,c}.
\]
The product is closed, because both factors are compact and \(C_0\)
is normal. Every preimage can be corrected by an element of
\(P_{4,c}\) into \(C_0\), proving the equality in both directions.
The established nonmembership statements prove
\eqref{eq:B176-preimage}. The same algebra proves the corresponding
ordinary preimage formula and ordinary nonmembership.

B.164's actual connecting map
\[
 \partial:(\mathcal K_a/\mathcal N)^\delta
 \longrightarrow H^1(\langle\delta\rangle,\mathcal N),
 \qquad t\mathcal N\longmapsto[t^{-1}\delta(t)]
\]
is continuous and injective with closed image, because
\(\mathcal K_a^\delta=1\). B.169 places \(\mathcal I_5\) in the
known source kernel \(\mathcal I\), both inside its domain. Applying
injectivity and compactness proves \eqref{eq:B176-cohomology} with
witness \(\partial(w)\). The full local question is still
\[
 \ker\bigl[H^1(\langle\delta\rangle,\mathcal N)
       \to H^1(\langle\delta\rangle,\mathcal K_a)\bigr]
                         =\partial(\mathcal I).
\]
The new witness already lies on its right side. It excludes replacing
the entire known source by the three \(m\)-only families, and does not
establish or refute that surjectivity equality or Main B.
\end{proof}

\begingroup\footnotesize\raggedright
This is the complete previously reviewed argument in
\path{work_CE/four_generator_contact_closure_obstruction38.txt}.
It uses B.175's original marked finite-refinement construction and its
retained 515-check certificate; it requires no new computation or numerical
construction of the final finite ambient separator. No full induced
topology for \(J_4^{\rm ord}\), \(\mathcal K_4\), or an infinite-index
source kernel is asserted.
\par\endgroup
